% CODEX REPORT::LEAD::HASHLEN-0001::2026-10-07
\documentclass[11pt]{article}
\usepackage[T1]{fontenc}
\usepackage[margin=0.85in]{geometry}
\usepackage{amsmath,amssymb,amsthm,booktabs,graphicx,xurl,hyperref}
\setlength{\emergencystretch}{3em}
\newtheorem{lemma}{Lemma}
\title{HASHLEN-0001: Exact Fingerprint Length, Adaptive Replay, and the Remaining WRAPPER Cover}
\author{Lead synthesis report}
\date{7 October 2026}
\begin{document}
\maketitle

\paragraph{Status.}
P versus NP remains \textsc{Open}.  HASH-INDEX-NO-SHRINK1,
ADAPTIVE-EXACT1, the bucket ledger and SUMMARY-COVER1 are
\textsc{Proved/Exact}.  CACHE-FINGERPRINT1, the chosen-collision cache tests
and the arithmetic simulations are \textsc{Computational}.  The general
source-dependent witness-preserving cover is \textsc{Open} and load-bearing.

\section{Sources actually reviewed}
The run reconciled the current GitHub \texttt{main4.tex},
\texttt{main5.tex}, \texttt{main6.tex}, the combined source-replay solver,
the current proof-spine protocol, and the recovered October 6--7 wrapper
artifacts.  The Library precommit test accepted every one of 480 wrong-token
claims under the native existential check, while the ratchet-link test rejected
all wrong continuations for an injective commitment and accepted all 1,000
under the nonbinding commitment.  The shared BAD-SEED construction preserved
exact replay with linear seed/history description while defeating route-count
as a generic lower bound.

The Drive recovery included WORKSAT-0011, WORKSAT-0017 and WORKSAT-0019.
Their rolling fingerprint reduced finite measured storage but provided only
fingerprint-weak sameness; the truncated sixteen-bit content hash failed
strong sameness; and the digest-plus-adaptive-pointer proposal either answered
equality only, re-solved on a miss, or rematerialized linear content for full
access.  The dynamic-ratchet notes likewise mark transfer from cell-probe or
online information arguments to unrestricted SAT as open.

Primary literature reviewed during the run included Carter--Wegman universal
hashing, Karp--Rabin verified fingerprints, Fredman--Koml\'os--Szemer\'edi
perfect hashing, Bryant canonical OBDDs, Bender et al. adaptive filters,
Bacchus--Dalmao--Pitassi component caching, Capelli et al. canonical TDDs, and
Biere--Lagniez--Lonca semantic hashing with provenance verification.  The last
paper is the closest exact architectural match: a cheap signature proposes a
decision-DNNF merge and two source-backed entailment tests certify it.  The
initial compilation may be exponential, so it does not close the acquisition
arrow for arbitrary CNF.

\section{Exact separation between indexing and merging}
Let $R$ be the number of classes of the future-use congruence among $Q$
generated states.  A complete exact verifier leaves one representative per
class, regardless of fingerprint length.  A sound incomplete verifier can
leave more.  Conversely, merging from an $h$-bit fingerprint alone requires
the implication
\[
 H(x)=H(y)\Longrightarrow x\equiv y,
\]
and therefore $2^h\ge R$.  Constructing a cheap small-range map with that
implication is the missing semantic quotient theorem rather than a way to
derive it.

For bucket sizes $m_z$, exact chained comparison work is
\[
 P(H)=\sum_z\binom{m_z}{2},
\qquad
 P(H)\ge\max\left\{0,{Q^2/2^h-Q\over2}\right\}.
\]
A fixed public map admits $P(H)=\binom Q2$.  For independent uniform hashes,
$\mathbb E[P]=\binom Q2/2^h$.  If $Q=M^u$, one replay costs $M^v$, and
$h=\kappa\log_2M$, the expected implementation exponent is at least
\[
 \max\{u+v,2u-\kappa+v\}.
\]
At the target $u=1/3,v=2/3$, a constant-length digest leaves exponent $4/3$;
$(1/3)\log_2M$ bits make expected replay linear; and
$(2/3)\log_2M$ bits make the expected number of false pairs constant.  These
are randomized expectations.  The fixed public worst case still gives joint
time-space exponent at least two in this cache model.

\begin{lemma}[Source-summary cover guard]
For every unique-witness unit CNF $F_x$, suppose an $h$-bit summary selects at
most $q$ restrictions, each leaving at most $\ell$ variables free.  Exact
coverage on all $x\in\{0,1\}^n$ requires
\[
 h+\log_2q+\ell\ge n.
\]
\end{lemma}
\begin{proof}
One summary bucket contains at least $2^{n-h}$ witnesses.  Its common $q$
restrictions cover at most $q2^\ell$ assignments.
\end{proof}
The same lemma applies when $h$ is the length of a source-independent random
seed.  Full-source adaptation escapes it, returning the problem to charged
source-dependent acquisition.

\section{Built and falsified variants}
The exact multilevel cache starts from a short prefix, extends a colliding
entry only when needed, and performs full normalized-residual replay before
every reuse.  Its correctness is independent of digest collision resistance.
The final guarded run completed twelve answers in 2.90 seconds.

\begin{center}
\resizebox{\textwidth}{!}{%
\begin{tabular}{lrrrrrr}
\toprule
scheme & max cache bits & bucket peak & equality replays & rejected & literal reads & peak work bits\\
\midrule
fixed 2 hex & 49,244 & 10 & 1,278 & 1,157 & 887,763 & 51,721\\
fixed 4 hex & 55,396 & 2 & 122 & 1 & 833,920 & 57,577\\
fixed 6 hex & 61,548 & 1 & 121 & 0 & 833,889 & 63,433\\
fixed 8 hex & 67,700 & 1 & 121 & 0 & 833,889 & 69,289\\
adaptive $2\to8$ & 55,116 & $10^*$ & 312 & 191 & 989,748 & 57,273\\
adaptive $4\to8$ & 55,412 & $2^*$ & 122 & 1 & 834,328 & 57,593\\
\bottomrule
\end{tabular}}
\end{center}
Four hexadecimal digits were the best fixed choice on these two workloads.
Six removed one rejected collision but saved only 31 literal reads while
adding 6,152 cache bits.  Eight added another 6,152 bits without saving work.
The two-digit adaptive cache reduced comparisons but paid 712 extension
rehashes, raising literal reads by 11.5 percent.

The collision-focused control produced hash-only answer errors at every
tested width: 1,288 at sixteen bits, 404 at twenty-four, and 72 at thirty-two
on its four-pass workloads.  Verified-hash and full-key caches remained
exact.  Equal-current-answer collision pairs were also separated by a later
common restriction, directly falsifying current SAT-bit equality as a future
congruence.

\section{Recovered components and final gap}
The recovered prime, hash, precommitment, BAD-SEED, source-selector and ratchet
constructions are not wasted.  They supply exact lineage, compact route
indices, binding tests, replay, and collision-safe dictionaries.  STREAMCOMP1,
ORBITBATCH1 and BATCHSETUP1 give conditional ledgers once a cover exists.
MACRO-LINEAGE-SPACE1 makes route bookkeeping logarithmic after contraction.
CLAUSE-DOM-SELECT1 is a genuine cheap merge, and bounded-width frontier/TDD
methods are genuine positive subclasses.

The strongest current composition is therefore
\[
\begin{gathered}
\text{full-source cheap witness-preserving cover}\\
\Longrightarrow q\le M^{1/3+o(1)},\quad |G_i|\le M^{2/3+o(1)}\\
\Longrightarrow \text{exact replay and charged recursion}\\
\Longrightarrow (C,\sigma)=(1,2/3)\Longrightarrow C+\sigma=5/3<2\cos(\pi/7).
\end{gathered}
\]
The first arrow is still open.  Hash length affects the third arrow after the
cover is known; it does not prove the first.  The detailed proof is now in
\texttt{main7.tex}; \texttt{main4.tex}, \texttt{main5.tex}, and
\texttt{main6.tex} have been reconciled without promoting any open claim.

\section{Reproducibility and handoff}
The executable artifacts are
\path{cache_fingerprint_experiment.py} with
\path{cache_fingerprint_results.json}, and
\path{hash_length_redteam_2026-10-07.py} with its JSON output.  Every
Python run uses an internal 25-second deadline.  The cache batch completed in
2.90 seconds; the arithmetic audit checked 1,872 convex bucket inequalities
and 48 common-signature bounds with zero errors.  The next useful target is a
source-dependent canonical-class theorem for one nontrivial general-SAT
frontier, with discovery time, residual size, total fanout, verifier calls and
workspace all charged.  Fingerprint optimization should then be rerun on the
proved class set.

% END CODEX REPORT::LEAD::HASHLEN-0001::2026-10-07
\end{document}
